\documentclass[11pt]{article}
\usepackage[margin=0.85in]{geometry}
\usepackage{amsmath,amssymb,amsthm}
\usepackage[hidelinks]{hyperref}
\newtheorem{theorem}{Theorem}
\newcommand{\R}{\mathbb R}
\title{A negative H-eigenvalue of a reflected Steiner distance tensor}
\author{Research manuscript}
\date{29 October 2026}
\begin{document}
\maketitle
\begin{abstract}
For the star $K_{1,29}$, the order-four Steiner distance tensor reflected in its last slot has a negative real H-eigenvalue. An exact two-orbit calculation reduces all thirty coordinate equations to a polynomial root in $(5,8)$, whose existence follows from the intermediate value theorem. This answers negatively the reflected-tensor positivity question of Cooper--Du, arXiv:2505.10501v1, pp.~9--10. It does not decide their hyperdeterminant-sign conjecture. The accompanying Lean certificate uses the actual minimum edge count over connected subgraphs and proves an unconditional negative H-pair.
\end{abstract}

\section{Objects and scope}
Let $T=K_{1,29}$ have vertex set $V=\{0,1,\ldots,29\}$, centre $0$, and edges $\{0,a\}$ for $1\leq a\leq29$. It is connected and acyclic. For an ordered quadruple, including repetitions, let $M_{ijkl}$ be the minimum number of edges of a connected subgraph containing its terminals. Write $M(u,v,w,z)=\sum_{i,j,k,l}M_{ijkl}u_iv_jw_kz_l$. Let $J$ be the all-ones matrix and
\[
 A=\frac{2J}{30}-I,\qquad C_1(u,v,w,z)=M(u,v,w,Az).
\]
Thus $(C_1)_{ijkl}=\sum_hM_{ijkh}A_{hl}$. We use the first-free-index H-equations
\begin{equation}\label{H}
 g_i(x):=C_1(e_i,x,x,x)=\lambda x_i^3\quad(i\in V),\qquad x\in\R^{30}\setminus\{0\}.
\end{equation}
These are the literal last-slot reflection of the source and the coordinate convention of Qi~\cite{Qi}. $C_1$ need not be fully symmetric; no implication from the sign of its diagonal quartic to \eqref{H} is assumed.

\begin{theorem}\label{main}
There exist $s\in(5,8)$ and $\lambda<0$ such that $x=(s+29,-1,\ldots,-1)$ satisfies all equations \eqref{H} for the above tensor on $T$.
\end{theorem}

One counterexample at order four refutes universal positivity over trees, and hence any assertion of positivity for all even orders over all trees. It asserts neither failure at every even order nor minimality of the vertex count. The source's separate determinant-sign conjecture and its positivity result for $C_1+C_3$ are not contradicted.

\section{Edge cuts and the thirty equations}
Deleting edge $\{0,a\}$ separates the singleton leaf $\{a\}$ from its complement. Any connected subgraph containing terminals on both sides must contain this edge. Conversely, the union of the unique paths between terminals is connected and contains exactly the edges whose cuts are crossed. For an all-equal tuple it is just one vertex and has zero edges. This proves, also with repeated terminals, that the minimum distance counts the crossed cuts.

For a vector $u$ write $S_u=\sum_i u_i$. Distributing finite sums over ordered tuples therefore gives the mixed identity
\begin{equation}\label{cut}
 M(u,v,w,z)=\sum_{a=1}^{29}\left[S_uS_vS_wS_z-u_av_aw_az_a
 -(S_u-u_a)(S_v-v_a)(S_w-w_a)(S_z-z_a)\right].
\end{equation}
The two subtractions remove exactly the tuples lying entirely on one side of a cut. This classical tree edge-cut argument is the mechanism of Cooper--Du's Theorem~2.1~\cite{CD}; its mixed form is supplied here explicitly.

For $x=(s+29,-1,\ldots,-1)$ put $y=Ax$. Then $S_x=S_y=s$, $y_a=s/15+1$ at each leaf, and the complementary sums are $s+1$ and $14s/15-1$. Substituting $u=e_i$, $v=w=x$, $z=y$ into \eqref{cut}, a cut at a leaf different from $i$ contributes
\[
 D=s^3-(s+1)^2(14s/15-1).
\]
When $i$ itself is a leaf, its own cut contributes $E=s^3-s/15-1$. There are twenty-nine cuts for the centre and twenty-eight other cuts for a fixed leaf. Thus, with
\[
 d=2s^3-26s^2+32s+30,\qquad
 L=86s^3-728s^2+894s+810,
\]
expansion gives \emph{every} coordinate:
\begin{equation}\label{orbits}
 g_0=29D=\frac{29d}{30},\qquad
 g_i=28D+E=\frac{L}{30}\quad(1\leq i\leq29).
\end{equation}
Leaf permutation symmetry accounts for all twenty-nine leaf equations, without omitting exceptional coordinates.

\section{Exact compatibility and negativity}
Define the real polynomial
\[
 P(s)=29d(s)+(s+29)^3L(s).
\]
Exact integer arithmetic yields
\[
 P(5)=-85\,295\,770<0,\qquad P(8)=273\,617\,240>0.
\]
Polynomial continuity and the classical intermediate value theorem supply $s\in(5,8)$ with $P(s)=0$. On $[5,8]$,
\[
 d(s)=30-48s+2s(s-5)(s-8)\leq30-48s\leq-210.
\]
Indeed the last product is nonpositive. In particular, for the selected root,
\[
 \lambda=\frac{29d(s)}{30(s+29)^3}<0.
\]
The centre equation follows from this definition and \eqref{orbits}. Since $P(s)=0$,
\[
 g_i=\frac{L(s)}{30}=-\frac{29d(s)}{30(s+29)^3}
      =-\lambda=\lambda(-1)^3\quad(i>0).
\]
The vector is nonzero, so this proves Theorem~\ref{main}. Neither a numerical eigenvalue nor uniqueness of the root is needed.

For completeness, diagonal positive definiteness fails directly as well. At $z=(34,-1,\ldots,-1)$, corresponding to $s=5$, \eqref{orbits} gives $g_0=-203$ and $g_i=-217/3$ for $i>0$, whence
\[
 C_1(z,z,z,z)=\sum_i z_i g_i=-14413/3<0.
\]
This sign is not the proof of the literal nonsymmetric H-pair. Qi's positivity equivalence applies to symmetric tensors~\cite{Qi}. If one instead symmetrizes the tensor while preserving its diagonal quartic $F$, that interpretation also fails: minimize $F$ on the compact sphere $\sum_i u_i^4=1$. The displayed negative direction makes the minimum negative. The constraint gradient is nonzero there; Lagrange multipliers and Euler's degree-four identity give a negative H-eigenvalue for the symmetric tensor. This is a separate classical observation, not an additional claimed outcome or a formal endpoint of this packet.

\section{Certificate, provenance, and limits}
The source-only certificate contains six mathematical modules, audited over all 118 actual declarations (including private and generated constants). \texttt{StarMinimum} proves the star is a tree and identifies required edges. \texttt{SourceMinimum} defines \texttt{sourceDistance} as a natural-number minimum over actual connected \texttt{star.Subgraph} objects containing every terminal. Its edge-label bijection proves equality with the undirected \texttt{edgeSet.ncard}. The theorem \texttt{source\_minimum\_attained} supplies a connected minimizer and comparison with every connected competitor. The equality bridge to required-edge count then justifies the tensor used by \texttt{CutTensor} and \texttt{Contraction}. \texttt{Algebra} proves the interval-root witness; \texttt{Endpoint} joins these results.

The endpoint \texttt{concrete\_negative\_H\_pair} has no graph, contraction, or root-existence premise. The endpoint \texttt{source\_refutation} states that this star is a tree and that its specialized predicate \texttt{StarHPositivity} is false. It is not a general all-$n$ Steiner-distance library. The audit prints actual elaborated types, module origins, transitive axiom closures, and unsafe/partial/axiom flags, crosschecking both environment and module inventories. Only \texttt{propext}, \texttt{Classical.choice}, and \texttt{Quot.sound} occur; all three declaration flags are false. The runner pins Lean 4.34.1 and nine package revisions, uses serial strict builds, and checks reused upstream source/cache hashes. It does not rebuild upstream dependencies. Details and resource receipts are in the packet.

The authenticated source is arXiv v1~\cite{CD}; its submission date is 15 May 2025, while the PDF internally bears 13 June 2025. No journal counterpart was located in the inherited preliminary review, so no version-of-record comparison is claimed. Independent final review and root source-only replay passed with bounded source and prior-art coverage. No matching prior resolution was found in that screen; no global priority is claimed. This is one campaign outcome, with no separate credit for tools or formalization.

The strongest substantive objection is retained: the source already provides the cut mechanism, the star has only two vertex orbits, and the remaining algebra is elementary; moreover the positivity question is ancillary to the source's determinant strategy. Formal cost establishes neither novelty nor mathematical depth. The candidate content is the compatible negative pair for the literal tensor, not the classical cut or intermediate-value tools.

\begin{thebibliography}{9}
\bibitem{CD} J. Cooper and Z. Du, \emph{Determinants of Steiner Distance Hypermatrices}, arXiv:2505.10501v1, 2025, Theorem 2.1 and pp.~9--10. \url{https://arxiv.org/abs/2505.10501v1}.
\bibitem{Qi} L. Qi, \emph{Eigenvalues of a real supersymmetric tensor}, J. Symbolic Comput. \textbf{40} (2005), 1302--1324, especially pp.~1303--1304. DOI: 10.1016/j.jsc.2005.05.007.
\end{thebibliography}
\end{document}
